% Part: first-order-logic
% Chapter: lindstrom
% Section: lindstrom-proof

\documentclass[../../../include/open-logic-section]{subfiles}

\begin{document}

\olfileid{mod}{lin}{prf}

\olsection{Lindstr\"om 定理}

\begin{lem}
\ollabel{lem:lindstrom}
设 $!E \in L(\Lang{L})$，其中 $\Lang{L}$ 有限，并假设存在 $n \in \Nat$，
使得对任意两个结构 $\Struct{M}$ 和~$\Struct{N}$，若 $\Struct{M} \equiv_n
\Struct{N}$ 且 $\Struct{M} \models_L !E$，
则也有 $\Struct{N} \models_L !E$。那么 $!E$ 等价于一个一阶语句，
即存在一阶的 $!D$ 使得 $\Mod(L){!E} = \Mod(L){!D}$。
\end{lem}

\begin{proof}
取定这样的 $n$，使得任意两个 $n$ 等价的结构 $\Struct{M}$ 和 $\Struct{N}$ 在~$!E$ 上取值一致。
回顾 \olref[bas][pis]{prop:qr-finite}：在有限语言中，量词秩不大于~$n$ 的一阶语句在逻辑等价意义下只有有限多个。
现在，对每个固定的结构~$\Struct{M}$，令 $!D_{\Struct{M}}$ 为所有在~$\Struct{M}$ 中为真
且 $\QuantRank{!E} \le n$ 的一阶语句~$!E$ 的合取（该合取是有限的），
从而 $\Struct{N} \models !D_{\Struct{M}}$ 当且仅当 $\Struct{N}
\equiv_n \Struct{M}$。然后令 $!D = \textstyle\bigvee
\Setabs{!D_{\Struct{M}}}{\Struct{M} \models_L !E}$；该析取在逻辑等价意义下也是有限的。

结论 $\Mod(L){!E} = \Mod(L){!D}$ 成立。事实上，如果 $\Struct{N} \models_L !D$，
则存在某个 $\Struct{M} \models_L !E$ 使得 $\Struct{N} \models !D_{\Struct{M}}$，
因此也有 $\Struct{N} \models_L !E$（根据引理假设）。
反之，若 $\Struct{N} \models_L !E$，则 $!D_{\Struct{N}}$ 是 $!D$ 的一个析取支，
又因为 $\Struct{N}
\models !D_{\Struct{N}}$，故也有 $\Struct{N} \models_L !D$。
\end{proof}

\begin{thm}[Lindstr\"om 定理]
  \ollabel{thm:lindstrom} 设 $\tuple{L, \models_L}$ 具有紧致性和 Löwenheim--Skolem 性质。
  则 $\tuple{L, \models_L} \le \tuple{F, \models}$（因此 $\tuple{L, \models_L}$ 等价于一阶逻辑）。
\end{thm}

\begin{proof}
根据 \olref{lem:lindstrom}，只需证明：对任意有限语言 $\Lang{L}$ 中的 $!E
\in L(\Lang{L})$，都存在 $n \in \Nat$，使得对任意两个结构 $\Struct{M}$ 和~$\Struct{N}$：若
$\Struct{M} \equiv_n \Struct{N}$ 则 $\Struct{M}$ 和 $\Struct{N}$
在~$!E$ 上取值一致。因为若如此，$!E$ 便等价于某个一阶语句，由此即可得
$\tuple{L, \models_L} \le \tuple{F, \models}$。由于我们在一个有限、纯关系的语言中工作，根据 \olref[bas][pis]{thm:b-n-f}，我们可以将
$\Struct{M} \equiv_n \Struct{N}$ 这一陈述替换为等价的代数陈述 $I_n(\emptyset,\emptyset)$。

给定 $!E$，反设对每个 $n$，都存在结构 $\Struct{M}_n$ 和 $\Struct{N}_n$ 使得
$I_n(\emptyset, \emptyset)$，但（不妨设）$\Struct{M}_n \models_L !E$
而 $\Struct{N}_n \not\models_L !E$。根据同构性质，我们可以假定所有 $\Struct{M}_n$ 对语言中常项的解释相同；进一步，由于语言中只有有限多个原子语句，我们还可以假定它们满足相同的原子语句（否则可取 $\Struct{M}$ 序列的一个子序列）。令 $\Struct{M}$ 为所有 $\Struct{M}_n$ 的并，即包含每个 $\Struct{M}_n$ 作为子结构的唯一最小结构。如 \olref[lsp]{thm:abstract-p-isom} 的证明所示，令 $\Struct{M}^*$ 为 $\Struct{M}$ 的扩张，其论域为
$\Domain{M} \cup
\Domain{M}^{<\omega}$，所在的语言是扩充了拼接谓词 $P$ 与~$Q$ 的语言。

类似地，定义 $\Struct{N}_n$、$\Struct{N}$ 和 $\Struct{N}^*$。现在令 $\Struct{M}$ 为这样一个结构：其论域包含 $\Struct{M}^*$ 和 $\Struct{N}^*$ 的论域、自然数集~$\Nat$ 及其自然序~$\le$，语言中增加了表示论域
$\Domain{M}$、$\Domain{N}$、$\Domain{M}^{<\omega}$ 和
$\Domain{N}^{<\omega}$ 的谓词，以及按下述意义编码 $\Struct{M}_n$ 和 $\Struct{N}_n$ 的论域的谓词：
\begin{align*}
  \Domain{M_n} & = \Setabs{a \in \Domain{M}}{R(a, n)}; &
  \Domain{N_n} & = \Setabs{a \in \Domain{N}}{S(a,n)}; \\
  \Domain{M}^{<\omega}_n & = \Setabs{a \in \Domain{M}^{<\omega}}{R(a,n)}; &
  \Domain{N}^{<\omega}_n & = \Setabs{a \in \Domain{N}^{<\omega}}{S(a,n)}.
\end{align*}
该结构~$\Struct{M}$ 还有一个三元关系 $J$，使得
$J(n, \mathbf{a}, \mathbf{b})$ 成立当且仅当
$I_n(\mathbf{a}, \mathbf{b})$。

现在，在语言~$\Lang{L}$ 加上 $R$、$S$、$J$ 等所得的语言中，存在一个语句~$!D$，它断言：$\le$ 是一个有首元但无末元的离散线性序，且 $\Struct{M}_n \models
!E$，$\Struct{N}_n \not\models !E$，并且对于该序中的每个 $n$，
$J(n, \mathbf{a}, \mathbf{b})$ 成立当且仅当
$I_n(\mathbf{a}, \mathbf{b})$。

利用紧致性性质，我们可以找到 $!D$ 的一个模型 $\Struct{M}^*$，其中该线性序包含一个非标准元素~$n^*$。特别地，此时 $\Struct{M^*}$ 将包含子结构
$\Struct{M_{n^*}}$ 和 $\Struct{N_{n^*}}$，使得 $\Struct{M_{n^*}}
\models_L !E$ 且 $\Struct{N_{n^*}} \not\models_L !E$。但现在我们可以定义一个由来自
$\Domain{M_{n^*}}$ 和 $\Domain{N_{n^*}}$ 的 $k$ 元组构成的序对集合 $\mathcal{I}$，规定
$\tuple{\mathbf{a}, \mathbf{b}} \in \mathcal{I}$ 当且仅当
$J(n^*-k, \mathbf{a}, \mathbf{b})$，其中 $k$ 是
$\mathbf{a}$ 和 $\mathbf{b}$ 的长度。由于 $n^*$ 是非标准的，对每个标准的 $k$，都有 $n^* - k >0$，并且集合 $\mathcal{I}$ 见证了
$\Struct{M_{n^*}} \simeq_p
\Struct{N_{n^*}}$。但根据 \olref[lsp]{thm:abstract-p-isom}，
$\Struct{M_{n^*}}$ 与 $\Struct{N_{n^*}}$ 是 $L$ 等价的，这就产生了矛盾。
\end{proof}

\end{document}